\section{Related Work}
\subsection{Structured Product and Hybrid Retrieval}
Product search must match queries to heterogeneous catalog fields. Lexical retrieval such as BM25 provides a strong order-insensitive baseline under a fixed token multiset~\cite{robertson2009bm25}; dense matching and field-aware models address semantic mismatch~\cite{choi2020semantic}. Attribute extraction can supply explicit matching signals~\cite{loughnane2024attributes}. CHARM constructs hierarchical field representations with block-triangular attention and uses aggregated vectors for candidate retrieval followed by richer field-level matching~\cite{freymuth2025charm}. Its learned, field-aware hierarchy differs from our frozen-encoder views of complete, fact-equivalent product records. We use it to position multi-vector product retrieval, not as an implemented baseline.

Hybrid retrieval combines lexical and semantic signals. Reciprocal rank fusion combines rankings without requiring comparable score scales~\cite{cormack2009rrf}. We include fixed full-list RRF as a practical competitor because a representation strategy is less compelling if a simple hybrid already obtains better recall. Our comparisons hold candidate count fixed and separately report the extra storage and scoring cost of multiple product vectors. The purpose is to evaluate the representation decision, rather than propose another fusion algorithm.

\subsection{Content and Visibility Optimization}
GEO studies how content changes influence visibility in generative systems~\cite{aggarwal2024geo}; E-GEO specializes this objective to e-commerce~\cite{bagga2025egeo}, while SAGEO Arena evaluates search-augmented optimization in a broader retrieval-and-generation setting~\cite{kim2026sageo}. Ranking-incentivized content modification likewise studies strategic changes under quality constraints~\cite{goren2020ranking}. Our initial question is different: why should a semantically irrelevant representation choice alter visibility in the first place? Our subsequent strategy evaluation preserves raw facts and uses query-independent rules or platform indexing. We do not evaluate generated product claims, agent recommendations, purchases, or merchant welfare.

\subsection{Representation Robustness}
Serialization sensitivity is not a new discovery in neural retrieval. Bhandari et al. study alternative table serializations, centroid representations, and a residual adapter toward centroid targets~\cite{bhandari2026tabular}. Their centroids motivate our full-record averaging baseline, which is distinct from independent attribute encoding followed by set pooling. We extend the empirical question to judged-relevant product inclusion, target-only causal control, persistent omission, cross-encoder candidate limitations, and recall--cost tradeoffs. We do not claim that a finite mean over views is a new invariant-learning algorithm.

Order effects and behavioral retrieval testing also caution against treating aggregate effectiveness as a complete description of a neural system~\cite{abdou2022wordorder,macavaney2022abnirml}. Our task-defined equivalence class is deliberately narrower than arbitrary rewriting: complete raw attribute atoms are moved while their contents and multiplicities remain fixed. Generalization to other structured Web domains remains an empirical question.
